%-------------------------------------------------------------------------------
% Brandon Jahoor — Resume
% Built on Jake Gutierrez's resume template (MIT):
%   https://github.com/jakegut/resume
% Build: latexmk -pdf brandon_jahoor_resume.tex
%-------------------------------------------------------------------------------
\documentclass[letterpaper,10pt]{article}

% Computer Modern ships fixed design sizes topping out at 24.88pt (\Huge).
% Ask for more and LaTeX silently substitutes 24.88 instead. fix-cm exposes the
% scalable versions, which is what lets the name be set to an arbitrary size.
\usepackage{fix-cm}
\usepackage{latexsym}
\usepackage[empty]{fullpage}
\usepackage{titlesec}
\usepackage{marvosym}
\usepackage{fontawesome5}
\usepackage[usenames,dvipsnames]{xcolor}
\usepackage{verbatim}
\usepackage{enumitem}
% Micro-adjusts letter spacing and lets glyphs protrude slightly into the
% margin. Invisible to a reader; recovers whole lines by pulling back words
% that would otherwise orphan onto a line of their own.
\usepackage[activate={true,nocompatibility},final,tracking=true,factor=1100,stretch=12,shrink=12]{microtype}
% University of Waterloo official palette
% https://uwaterloo.ca/brand/how-express-our-brand/colour
\definecolor{uwpurple}{HTML}{5D0096}   % Engineering L4 — the name only
\usepackage[hidelinks]{hyperref}
\usepackage{fancyhdr}
\usepackage[english]{babel}
\usepackage{tabularx}
\input{glyphtounicode}

\pagestyle{fancy}
\fancyhf{}
\fancyfoot{}
\renewcommand{\headrulewidth}{0pt}
\renewcommand{\footrulewidth}{0pt}

% Jake's margins, tightened to 0.35in sides so no bullet orphans a word
\addtolength{\oddsidemargin}{-0.65in}
\addtolength{\evensidemargin}{-0.65in}
\addtolength{\textwidth}{1.3in}
\addtolength{\topmargin}{-.79in}
\addtolength{\textheight}{1.39in}

\urlstyle{same}
\raggedbottom
\raggedright
\setlength{\tabcolsep}{0in}

% Section headings: small caps, full-width rule beneath
\titleformat{\section}{
  \vspace{-4pt}\scshape\raggedright\large
}{}{0em}{}[\color{black}\titlerule \vspace{-5pt}]
% Tighter than titlesec's default space above and below each heading
\titlespacing*{\section}{0pt}{6pt}{11pt}

% Machine-readable PDF so ATS keyword extraction works
\pdfgentounicode=1

% Project links: purple icon + purple underlined label
\newcommand{\ghlink}[1]{{\color{uwpurple}\faGithub\ \href{#1}{github}}}
\newcommand{\stlink}[1]{{\color{uwpurple}\faGlobe\ \href{#1}{site}}}
\newcommand{\gllink}[1]{{\color{uwpurple}\faGitlab\ \href{#1}{gitlab}}}
\newcommand{\pflink}[1]{{\color{uwpurple}\faIcon{external-link-alt}\ \href{#1}{portfolio}}}


%-------------------------------------------------------------------------------
% Jake's macros
%-------------------------------------------------------------------------------
% Jake writes this as {#1 \vspace{-2pt}} — with the \vspace INSIDE the text
% group. A bullet ending near the line edge then wraps that space onto an
% invisible second line, costing ~11pt of phantom height. Moving \vspace out of
% the group keeps it in vertical mode, where it cannot break.
\newcommand{\resumeItem}[1]{
  \item\small{#1}\vspace{-2pt}
}

% {Role}{Dates}{Organization}{Location}
\newcommand{\resumeSubheading}[4]{
  \vspace{-2pt}\item
    \begin{tabular*}{\textwidth}[t]{l@{\extracolsep{\fill}}r}
      \textbf{#1} & #2 \\
      \textit{\small#3} & \textit{\small #4} \\
    \end{tabular*}\vspace{-7pt}
}

% {Project name \emph{--} stack}{Right-hand note}
\newcommand{\resumeProjectHeading}[2]{
    \item
    \begin{tabular*}{\textwidth}{l@{\extracolsep{\fill}}r}
      \small#1 & #2 \\
    \end{tabular*}\vspace{-7pt}
}

% Defined by Jake's template; unused here but kept so the preamble matches.
\newcommand{\resumeSubSubheading}[2]{
    \item
    \begin{tabular*}{\textwidth}{l@{\extracolsep{\fill}}r}
      \textit{\small#1} & \textit{\small #2} \\
    \end{tabular*}\vspace{-7pt}
}

\newcommand{\resumeSubItem}[1]{\resumeItem{#1}\vspace{-4pt}}

\renewcommand\labelitemii{$\vcenter{\hbox{\tiny$\bullet$}}$}

\newcommand{\resumeSubHeadingListStart}{\begin{itemize}[leftmargin=0in, label={}]}
\newcommand{\resumeSubHeadingListEnd}{\end{itemize}}
\newcommand{\resumeItemListStart}{\begin{itemize}}
\newcommand{\resumeItemListEnd}{\end{itemize}\vspace{-5pt}}

%===============================================================================
\begin{document}

%---------- HEADING ----------
% Two top-aligned minipages rather than a tabular*: tabular cells share a
% baseline, which drops the small contact text ~10pt below the top of the
% \Huge name. The \vspace{0pt} is the idiom that forces true top alignment.
\noindent
\begin{minipage}[t]{0.56\textwidth}
  \vspace{0pt}
  \textbf{\fontsize{27.67}{27.67}\selectfont \scshape \color{uwpurple}Brandon Jahoor}\\[-1.2pt]
  % 10pt explicitly, not \small. This document is on the 10pt class where
  % \small is 9pt; Jake's template is on 11pt, where \small is 10pt. Setting
  % the size directly matches the template's tagline size without moving the
  % whole document to 11pt, which cannot be made to fit on one page.
  \textit{\fontsize{10}{12}\selectfont Robotics Software Engineering for Humanoid Embodiment of Physical AI}
\end{minipage}%
\hfill
\begin{minipage}[t]{0.42\textwidth}
  \vspace{0pt}
  % Four links in a 2x2 grid. Both columns are right-aligned, and the long
  % entries (LinkedIn, email) share the left column so the gap down the middle
  % stays even — pairing a long entry against a short one is what made the
  % previous order look ragged.
  % Pushed down so the two rows sit centred against the name-and-tagline block
  % opposite, rather than being stretched to touch its top and bottom.
  \vspace{6pt}
  \raggedleft \fontsize{8}{8}\selectfont
  \begin{tabular}{@{}r@{\hspace{6pt}}r@{}}
    {\color{uwpurple}\faLinkedin\ \href{https://linkedin.com/in/bjahoor}{\underline{linkedin.com/in/bjahoor}}} &
    {\color{uwpurple}\faEnvelope\ \href{mailto:bjahoor@uwaterloo.ca}{\underline{bjahoor@uwaterloo.ca}}} \\[5pt]
    {\color{uwpurple}\faGlobe\ \href{https://bjahoor.github.io}{\underline{bjahoor.github.io}}} &
    {\color{uwpurple}\faGithub\ \href{https://github.com/bjahoor}{\underline{github.com/bjahoor}}} \\
  \end{tabular}
\end{minipage}

%---------- EDUCATION ----------
% No section heading, but keep the rule so the degree line still reads as its
% own band rather than running into the contact block.
\vspace{3.6pt}
\hrule height 0.8pt
\vspace{5pt}
{\small \textbf{University of Waterloo} $|$ \textit{Candidate for BASc in Mechatronics
Engineering Honours, \href{https://uwaterloo.ca/academic-calendar/undergraduate-studies/catalog\#/programs/Hye0Jk00ih?group=Options\%3A\%20Faculty\%20of\%20Engineering&bc=true&bcCurrent=Artificial\%20Intelligence\%20Option&bcGroup=Options\%3A\%20Faculty\%20of\%20Engineering&bcItemType=programs}{Minor in Artificial Intelligence}} \hfill \textit{2028}}
\vspace{2pt}

%---------- SKILLS ----------
\section{Skills}
 \begin{itemize}[leftmargin=0in, label={}]
    \small{\item{
     \textbf{Robot Learning}{: VLAs ($\pi_{0.5}$, GR00T, SmolVLA), ACT, RL, imitation learning, cross-embodiment, sim-to-real, synthetic data} \\
     \textbf{Sim \& Robotics}{: Isaac Sim/Lab, MuJoCo, Genesis, cuRobo, ROS2, WBC, URDF/USD, IK/FK} \\
     \textbf{Software \& Infra}{: Python, C++, PyTorch, JAX, openpi, LeRobot, Docker, Linux, sharded training, Nvidia Jetson, Claude Code}
    }}
 \end{itemize}

%---------- EXPERIENCE ----------
\section{Experience}
  \resumeSubHeadingListStart

    \resumeSubheading
      {Robotics Software Engineering (Co-op)}{Sept 2026 -- April 2027}
      {Axibo AI $|$ Humanoids $|$ \stlink{https://www.axibo.com/} $\cdot$ \pflink{https://bjahoor.github.io/projects/axibo-physical-ai}}{Ontario}
      \resumeItemListStart
        \resumeItem{Fine-tune, evaluate and deploy SOTA VLAs (\href{https://www.physicalintelligence.company/blog/pi05}{\textbf{\boldmath$\pi_{0.5}$}}, \href{https://developer.nvidia.com/isaac/gr00t}{\textbf{GR00T N1.7}}, \href{https://www.dexmal.com/models/dm0.5}{\textbf{DM0.5}}) end-to-end as generalizable bimanual and whole-body policies, both integrated with the \textbf{WBC} and locomotion policies, with a focus on cross-embodiment and sim-to-real.}
        \resumeItem{Created and designed a batched synthetic data generation toolchain for delta-EE bimanual policies using \href{https://genesis-world.readthedocs.io/en/latest/}{\textbf{Genesis World}}, \href{https://developer.nvidia.com/isaac/sim}{\textbf{Isaac}}, \href{https://mujoco.org/}{\textbf{MuJoCo}} and \href{https://curobo.org/}{\textbf{cuRobo}} (motion planning). Features like augmentation, domain-randomization, noise improve \textbf{cross-embodiment}.}
        \resumeItem{Accelerated SOTA VLA research iteration and cut compute cost with custom training pipelines and scripts forked from SOTA repos, primarily \href{https://github.com/Physical-Intelligence/openpi}{\textbf{openpi}} (\textbf{JAX/Orbax}) (\textbf{2.5$\times$} faster steps) and \href{https://github.com/NVIDIA/Isaac-GR00T}{\textbf{Isaac-GR00T}} (\textbf{PyTorch}), without sacrificing benchmarks.}
        \resumeItem{Optimized VLA training through sharding across multi-datacenter GPUs, upgraded version pins, selective weight and optimizer-state quantization, partial backbone freezing/pruning, selective rematerialization and batched checkpointing.}
        \resumeItem{Built batched sim evaluation toolchains with custom success criteria (per-component scoring with penalties) for SOTA bimanual \& whole-body policies, ran \textbf{sim-to-sim} \& \textbf{sim-to-real} evaluation on a variety of humanoids \& bimanual arms, and raised benchmark scores by tuning training recipes \& inference settings.}
        \resumeItem{Optimizing real-time inference via \textbf{RTC} \& Mixture of Horizons (\textbf{MoH}), tuning denoising steps \& prediction horizons.}
        \resumeItem{Productionizing the custom \textbf{UMI} pipeline and developed training recipes to maximize the efficacy of public data and synthetically generated data (including re-simulating public datasets with custom modifications) while minimizing required real-world data.}
        \resumeItem{Architected robotic runtimes with \textbf{LLM/agent-harness} orchestration, either sequencing VLA skills or driving the whole robot directly. Built toolchains, pipelines and workflows across the stack using \textbf{agent-centric} development.}
      \resumeItemListEnd

    \resumeSubheading
      {Robotics Software Engineering (Co-op)}{Jan 2026 -- April 2026}
      {Thalassa Robotics $|$ Subsea Robotics $|$ \stlink{https://www.thalassainc.com/} $\cdot$ \pflink{https://bjahoor.github.io/projects/thalassa-simulation-toolchain}}{California}
      \resumeItemListStart
        \resumeItem{Built a high-fidelity toolchain on \textbf{Isaac Sim} used team-wide, integrating
          and configuring \textbf{optics and hydrodynamics} it lacks.}
        \resumeItem{Prepared the robot's \textbf{USD} model and joint articulation so the simulated
          assembly matched the real one.}
        \resumeItem{Layered control abstractions so simulated motors drive exactly like the real
          servos, with acceleration-bounded interpolation.}
      \resumeItemListEnd

    \resumeSubheading
      {Robotics Software Engineering (Research Assistant)}{Sept 2025 -- Dec 2025}
      {University of Waterloo RoboHub $|$ Robotics Research Lab $|$ \gllink{https://git.uwaterloo.ca/bjahoor/xsens-panda} $\cdot$ \stlink{https://uwaterloo.ca/robohub/} $\cdot$ \pflink{https://bjahoor.github.io/projects/uw-robohub-research}}{Ontario}
      \resumeItemListStart
        \resumeItem{Streamed \textbf{human arm motion} from \textbf{Xsens MOCAP} into \textbf{ROS2}
          in real time, for retargeting onto a 7-axis \textbf{Franka Emika Panda}.}
        \resumeItem{Containerised the \textbf{ROS2} stack in \textbf{Docker} so results
          reproduced across the lab's workstations.}
      \resumeItemListEnd

    \resumeSubheading
      {Robotics Engineering (Co-op)}{June 2025 -- Aug 2025}
      {University of Waterloo Robotics Team $|$ Mars Rover Team $|$ \ghlink{https://github.com/bjahoor/cam_package} $\cdot$ \stlink{https://uwroboticsmarketing.github.io/} $\cdot$ \pflink{https://bjahoor.github.io/projects/mars-style-rover}}{Ontario}
      \resumeItemListStart
        \resumeItem{Developed a \textbf{ROS2 RealSense RGB-D} pipeline and solved \textbf{Mars rover
          drivetrain kinematics} on an \textbf{Nvidia Jetson}.}
      \resumeItemListEnd

    \resumeSubheading
      {Robotics Engineering (Co-op)}{May 2025 -- June 2025}
      {BHF Robotics $|$ Agricultural Robotics $|$ \stlink{https://bhfrobotics.com/} $\cdot$ \pflink{https://bjahoor.github.io/projects/autonomous-electric-weeder}}{Ontario}
      \resumeItemListStart
        \resumeItem{Implemented \textbf{ROS2 autodrive} and trained a \textbf{deep-learning}
          weed/crop detector for an autonomous farming robot.}
      \resumeItemListEnd

  \resumeSubHeadingListEnd

%---------- PROJECTS ----------
\section{Projects}
    \resumeSubHeadingListStart

      \resumeProjectHeading
        {\textbf{VLA Finetuning} $|$ \emph{SmolVLA, LeRobot, SO-ARM101, Nvidia Jetson} $|$
          \ghlink{https://github.com/bjahoor/lerobot-cookbook} $\cdot$
          \pflink{https://bjahoor.github.io/projects/vla-finetuning}}{}
        \resumeItemListStart
          \resumeItem{Trained \textbf{ACT} and \textbf{SmolVLA} on 50--100 teleoperated demos/task, running
          \textbf{30 Hz} on a Jetson Orin Nano.}
        \resumeItem{Optimized batch and worker counts for \textbf{LeRobot PyTorch} training on an \textbf{RTX 3060 Ti}.}
        \resumeItemListEnd
      \resumeProjectHeading
        {\textbf{Chess Robot} $|$ \emph{SO-ARM101, SolidWorks, 3D printing, LeRobot, Placo, Python}
          $|$ \emph{Stealth Startup} $|$ \pflink{https://bjahoor.github.io/projects/chess-robot}}{}
        \resumeItemListStart
          \resumeItem{\textbf{CAD}-modified an \textbf{SO-ARM101} and built its rig, then developed a
            three-call Python API that plugs straight into a chess engine.}
          \resumeItem{Extended reach and fused the wrist to drop a motor, updating the \textbf{URDF}
            model, solver constraints and \textbf{LeRobot} driver.}
          \resumeItem{Built the motion engine over \textbf{LeRobot} and \textbf{Placo} IK:
            convergence loops, \textbf{1 mm} waypoints, \textbf{PID gravity compensation}.}
        \resumeItemListEnd
      \resumeProjectHeading
        {\textbf{Autonomous Navigation Stack} $|$ \emph{ROS2, C++, LiDAR, A*, Foxglove} $|$ \emph{WATonomous} $|$
          \ghlink{https://github.com/bjahoor/wato_asd_navigation} $\cdot$ \pflink{https://bjahoor.github.io/projects/watonomous-navigation}}{}
      \resumeProjectHeading
        {\textbf{Robot Arm Simulation} $|$ \emph{Isaac Sim, Isaac Lab, CUDA, RL} $|$ \emph{Stealth Startup} $|$
          \pflink{https://bjahoor.github.io/projects/robot-arm-simulation}}{}
      \resumeProjectHeading
        {\textbf{Hexapod Robot} $|$ \emph{Nvidia Jetson, OWL-ViT, SO-ARM101, RealSense} $|$ \pflink{https://bjahoor.github.io/projects/hexapod-robot}}{}
      \resumeProjectHeading
        {\textbf{AI Roommate} $|$ \emph{Nvidia Jetson, Qwen2.5-0.5B, YOLOv8s, FastAPI} $|$
          \ghlink{https://github.com/bjahoor/ai_roommate_assistant} $\cdot$ \pflink{https://bjahoor.github.io/projects/ai-roommate}}{}
      \resumeProjectHeading
        {\textbf{Autonomous Rover Simulation} $|$ \emph{Gazebo, ROS2, Nav2, Python} $|$ \emph{RoboHub} $|$
          \ghlink{https://github.com/bjahoor/ROS2_Construct_Rover} $\cdot$ \pflink{https://bjahoor.github.io/projects/autonomous-rover-simulation}}{}
      \resumeProjectHeading
        {\textbf{AI Object Detection} $|$ \emph{ROS2, OWL-ViT, RealSense, OpenCV, Python} $|$
          \ghlink{https://github.com/bjahoor/ai_cam_package} $\cdot$ \pflink{https://bjahoor.github.io/projects/ai-object-detection}}{}




    \resumeSubHeadingListEnd
%---------- CLOSING ----------
\vspace{-16pt}
\noindent\rule{\textwidth}{0.8pt}\par
\vspace{5pt}
{\centering
  {\rmfamily\large\bfseries\itshape See my portfolio:\
  \textcolor{uwpurple}{\href{https://bjahoor.github.io}{\underline{bjahoor.github.io}}}}\par}

\end{document}
